% ============================================================
%  From the Quadratic Sequence to the Conical Helix:
%  A Direct Algebraic Construction
%  Dogan Balban — 2026
% ============================================================

\documentclass[12pt,a4paper]{article}

% ── Encoding and Language ────────────────────────────────────
\usepackage[utf8]{inputenc}
\usepackage[T1]{fontenc}
\usepackage[english]{babel}
\usepackage{csquotes}

% ── Bibliography ─────────────────────────────────────────────
\usepackage[
  backend=biber,
  style=alphabetic,
]{biblatex}
\addbibresource{references.bib}

% ── Mathematics ──────────────────────────────────────────────
\usepackage{mathtools}
\usepackage{amsmath}
\usepackage{amssymb}
\usepackage{amsthm}

% ── Page Layout ──────────────────────────────────────────────
\usepackage{geometry}
\geometry{
  left   = 2.5cm,
  right  = 2.5cm,
  top    = 3.0cm,
  bottom = 3.0cm
}

% ── Hyperlinks ───────────────────────────────────────────────
\usepackage[
  colorlinks = true,
  linkcolor  = blue!60!black,
  citecolor  = blue!60!black,
  urlcolor   = blue
]{hyperref}

% ── Graphics and Colours ─────────────────────────────────────
\usepackage{graphicx}
\graphicspath{{bilder/}}
\usepackage[table]{xcolor}

% ── Boxes ────────────────────────────────────────────────────
\usepackage{tcolorbox}
\tcbuselibrary{skins, breakable}

% ── Tables ───────────────────────────────────────────────────
\usepackage{booktabs}
\usepackage{tabularx}
\usepackage{array}
\usepackage{float}

% ── Units and Numbers ────────────────────────────────────────
\usepackage{siunitx}
\sisetup{output-decimal-marker = {.}}

% ── Lists ────────────────────────────────────────────────────
\usepackage{enumitem}

% ── Miscellaneous ────────────────────────────────────────────
\usepackage{textcomp}
\usepackage{gensymb}

% ============================================================
%  THEOREM ENVIRONMENTS
% ============================================================

\theoremstyle{plain}
\newtheorem{theorem}{Theorem}[section]
\newtheorem{lemma}[theorem]{Lemma}
\newtheorem{proposition}[theorem]{Proposition}
\newtheorem{corollary}[theorem]{Corollary}

\theoremstyle{definition}
\newtheorem{definition}[theorem]{Definition}
\newtheorem{example}[theorem]{Example}

\theoremstyle{remark}
\newtheorem{remark}[theorem]{Remark}

% ============================================================
%  TITLE
% ============================================================

\title{%
  \textbf{From the Quadratic Sequence to the Conical Helix:\\[0.3em]
  A Direct Algebraic Construction}\\[0.8em]
  \large Via the Universal Quantity $Q(k) = \sqrt{48k+1}$
}
\author{Dogan Balban}
\date{2026}

% ============================================================
%  DOCUMENT
% ============================================================

\begin{document}

\maketitle

\begin{abstract}
We consider the quadratic sequence $k(n) = \frac{1}{6}(2n^2+3n+1)$,
which takes integer values precisely when $n \equiv 1, 5 \pmod{6}$ ---
the residue classes containing all primes greater than~$3$.
We show that the algebraic identity $(4n+3)^2 = 48\,k(n) + 1$ gives
rise to a universal quantity $Q(k) = \sqrt{48k+1}$ that encodes the
complete geometry of a conical helix: radius, angle, and height are
all explicit closed-form expressions in~$Q$.
The helix lies simultaneously on an Archimedean spiral
$r = \frac{6}{\pi^2}\theta$ and a cone $z = -\frac{3}{4} + \frac{\pi}{2}r$,
both as algebraic consequences of the single identity.
The arc length of the parabola $k(n)$ converges asymptotically to the
difference values $\Delta k$, and all primes $p > 3$ appear as
distinguished discrete points on the helix.
The construction is entirely deductive: every geometric object is
\emph{derived} from the sequence, not empirically observed.
\end{abstract}

\tableofcontents
\clearpage

% ── Chapter 1: Introduction ──────────────────────────────────
% ============================================================
%  Chapter 1 — Introduction
% ============================================================
\section{Introduction}\label{sec:introduction}

\subsection{Overview}

The geometric visualisation of number-theoretic structures has a long
tradition, from the Ulam spiral~\cite{Stein1964} to the Sacks
spiral~\cite{Sacks2003}. These classical constructions are primarily
\emph{empirical}: patterns are observed, not derived.

The present work takes a different approach. Starting from the quadratic
sequence
\begin{equation}\label{eq:k-def}
  k(n) = \frac{1}{6}(2n^2 + 3n + 1),
\end{equation}
which gives the arithmetic mean of the first $n$ square numbers, we show
that a single algebraic identity --- the \emph{fundamental identity}
$(4n+3)^2 = 48\,k(n)+1$ --- determines a conical helix whose every
geometric property is an exact consequence of the sequence.

No intermediate constructions (spiral winding, circumference identification,
cone envelope) are required. The entire geometry is encoded in the
\emph{universal quantity}
\begin{equation}\label{eq:Q-intro}
  Q(k) = \sqrt{48k+1}.
\end{equation}

\subsection{Main Result}

The central theorem of this paper can be stated as follows.

\begin{theorem}[Main Theorem]\label{thm:main}
Let $k(n) = \frac{1}{6}(2n^2+3n+1)$ and $Q(k) = \sqrt{48k+1}$. Then:
\begin{enumerate}[label=\textup{(\roman*)}]
  \item $k(n) \in \mathbb{Z}$ if and only if $n \equiv 1, 5 \pmod{6}$.
  \item For all such $n$, the point
  \begin{equation}\label{eq:helix-main}
    \vec{p}(k) =
    \begin{pmatrix}
      \dfrac{Q}{2\pi}\cos\!\Bigl(\dfrac{\pi Q}{12}\Bigr)\\[8pt]
      \dfrac{Q}{2\pi}\sin\!\Bigl(\dfrac{\pi Q}{12}\Bigr)\\[8pt]
      \dfrac{Q-3}{4}
    \end{pmatrix}
  \end{equation}
  lies on a conical helix that satisfies simultaneously:
  \begin{itemize}
    \item the Archimedean spiral law $r = \frac{6}{\pi^2}\,\theta$,
    \item the cone equation $z = -\frac{3}{4} + \frac{\pi}{2}\,r$.
  \end{itemize}
  \item All primes $p > 3$ appear as discrete points on this helix.
\end{enumerate}
\end{theorem}

The proof of this theorem is distributed across
Sections~\ref{sec:sequence}--\ref{sec:primes} and assembled from
self-contained intermediate results.

\subsection{Structure of the Paper}

Section~\ref{sec:sequence} establishes the arithmetic properties of the
sequence $k(n)$: integrality condition, difference structure, and the
connection to primes.
Section~\ref{sec:identity} derives the fundamental identity and defines
the universal quantity $Q(k)$, together with the exact inverse formula
$n(k)$.
Section~\ref{sec:helix} constructs the conical helix directly from $Q$
and proves the spiral and cone conformity as algebraic corollaries.
Section~\ref{sec:arclength} treats the arc length of the parabola $k(n)$
and proves the asymptotic equivalence $L \approx \Delta k$.
Section~\ref{sec:primes} derives the prime number corollary.

\subsection{Notation}

Throughout this paper, $k(n)$ denotes the quadratic mean-value sequence,
$Q(k) = \sqrt{48k+1}$ the universal quantity, $\theta$ the angular
parameter of the helix, $r$ the radius, and $z$ the height coordinate.
The first differences of the integer-valued subsequence alternate between
\emph{long steps} $\delta^{(4)}_m = 16m-6$ (index difference~$4$) and
\emph{short steps} $\delta^{(2)}_m = 8m+1$ (index difference~$2$); the
superscript indicates the step width $\Delta n$.


% ── Chapter 2: The Quadratic Sequence ────────────────────────
% ============================================================
%  Chapter 2 — The Quadratic Sequence
% ============================================================
\section{The Quadratic Sequence}\label{sec:sequence}

\subsection{Definition and Origin}

The sequence arises from the classical summation formula for squares.
The sum of the first $n$ squares is
$S_n = \sum_{i=1}^{n} i^2 = \frac{n(n+1)(2n+1)}{6}$,
and its arithmetic mean
\begin{equation}\label{eq:k-explicit}
  k(n) = \frac{S_n}{n} = \frac{(n+1)(2n+1)}{6}
       = \frac{2n^2 + 3n + 1}{6}
\end{equation}
is a quadratic polynomial in $n$ with leading coefficient $\frac{1}{3}$.

\begin{remark}[OEIS entry]
The integer values of $k(n)$ are listed in the OEIS under
\href{https://oeis.org/A164576}{A164576}~\cite{OEIS} as
\emph{``Integer averages of the set of the first positive squares up to
some $n^2$''}.
\end{remark}

\subsection{Integrality Condition}

\begin{theorem}[Congruence condition]\label{thm:congruence}
The sequence $k(n)$ takes integer values if and only if
\begin{equation}\label{eq:cong}
  n \equiv 1, 5 \pmod{6}.
\end{equation}
\end{theorem}

\begin{proof}
We examine $2n^2 + 3n + 1 \pmod{6}$ for all residue classes:

\begin{center}
\begin{tabular}{ccccc}
\toprule
$n \!\pmod{6}$ & $2n^2 \!\pmod{6}$ & $3n \!\pmod{6}$ & $2n^2\!+\!3n\!+\!1 \!\pmod{6}$ & $6 \mid ?$ \\
\midrule
0 & 0 & 0 & 1 & No \\
1 & 2 & 3 & 0 & Yes \\
2 & 2 & 0 & 3 & No \\
3 & 0 & 3 & 4 & No \\
4 & 2 & 0 & 3 & No \\
5 & 2 & 3 & 0 & Yes \\
\bottomrule
\end{tabular}
\end{center}
Divisibility by~$6$ occurs precisely for $n \equiv 1$ and $n \equiv 5$.
\end{proof}

The first integer values are:

\begin{center}
\begin{tabular}{c*{10}{c}}
\toprule
$n$ & 1 & 5 & 7 & 11 & 13 & 17 & 19 & 23 & 25 & 29 \\
\midrule
$k(n)$ & 1 & 11 & 20 & 46 & 63 & 105 & 130 & 188 & 221 & 295 \\
\bottomrule
\end{tabular}
\end{center}

\subsection{Difference Structure: Long and Short Steps}\label{subsec:steps}

The integer-valued indices follow a periodic pattern. Within each
period $m = 1, 2, 3, \ldots$, there is a step of width $\Delta n = 4$
from $n = 6m-5$ to $n = 6m-1$, followed by a step of width
$\Delta n = 2$ from $n = 6m-1$ to $n = 6m+1$.

\begin{definition}[Long and short steps]\label{def:steps}
For the $m$-th period ($m \geq 1$) we define:
\begin{align}
  \delta^{(4)}_m &\coloneqq k(6m-1) - k(6m-5)
  \quad\text{(\emph{long step},\; $\Delta n = 4$)},
  \label{eq:long}\\[4pt]
  \delta^{(2)}_m &\coloneqq k(6m+1) - k(6m-1)
  \quad\text{(\emph{short step},\; $\Delta n = 2$)}.
  \label{eq:short}
\end{align}
The superscript indicates the index difference $\Delta n$.
\end{definition}

\begin{theorem}[Step-size formulae]\label{thm:step-formulae}
For all $m \geq 1$:
\begin{equation}\label{eq:step-formulae}
  \delta^{(4)}_m = 16m - 6,
  \qquad
  \delta^{(2)}_m = 8m + 1.
\end{equation}
\end{theorem}

\begin{proof}
Using $k(n) = \frac{1}{6}(2n^2+3n+1)$ and the difference-of-squares identity:
\begin{align*}
  k(6m\!-\!1) - k(6m\!-\!5)
  &= \frac{1}{6}\bigl[2\cdot(12m\!-\!6)\cdot 4 + 3\cdot 4\bigr]
  = \frac{96m-36}{6} = 16m-6,\\[4pt]
  k(6m\!+\!1) - k(6m\!-\!1)
  &= \frac{1}{6}\bigl[2\cdot 12m\cdot 2 + 3\cdot 2\bigr]
  = \frac{48m+6}{6} = 8m+1.
  \qedhere
\end{align*}
\end{proof}

The first differences are thus:
\begin{center}
\begin{tabular}{c*{8}{c}}
\toprule
$\delta_k$ & 10 & 9 & 26 & 17 & 42 & 25 & 58 & 33 \\
$\Delta n$ & 4 & 2 & 4 & 2 & 4 & 2 & 4 & 2 \\
\bottomrule
\end{tabular}
\end{center}

\begin{remark}[Pairwise sums form an arithmetic progression]
Each long--short pair sums to
$U_m = \delta^{(4)}_m + \delta^{(2)}_m = (16m-6)+(8m+1) = 24m-5$,
with constant increment $\Delta U = 24$. This arithmetic structure is a
direct algebraic consequence of Theorem~\ref{thm:step-formulae} and will
reappear as the circumference of the $m$-th spiral winding in
Section~\ref{sec:helix}.
\end{remark}

\subsection{Prime Number Property}\label{subsec:primes-mod6}

\begin{theorem}[Primes and the congruence classes]\label{thm:primes-mod6}
All primes $p > 3$ satisfy $p \equiv 1$ or $5 \pmod{6}$.
\end{theorem}

\begin{proof}
Every prime $p > 3$ is coprime to $6$, since it is divisible neither by~$2$
nor by~$3$. The residue classes coprime to~$6$ are precisely $1$ and
$5$~\cite[Ch.\,1]{IrelandRosen1990}.
\end{proof}

Thus, the indices at which $k(n)$ is an integer are exactly the residue
classes that contain all primes greater than~$3$. This coincidence is the arithmetic foundation of the prime number
corollary in Section~\ref{sec:primes}.


% ── Chapter 3: The Fundamental Identity and Q(k) ─────────────
% ============================================================
%  Chapter 3 — The Fundamental Identity and Q(k)
% ============================================================
\section{The Fundamental Identity and the Universal Quantity \texorpdfstring{$Q(k)$}{Q(k)}}
\label{sec:identity}

This section establishes the algebraic core of the entire construction:
a single identity that connects the index $n$, the function value $k$,
and all geometric coordinates of the helix.

\subsection{The Identity}

\begin{theorem}[Fundamental Identity]\label{thm:fundamental}
For all $n \geq 1$:
\begin{equation}\label{eq:fundamental}
  (4n+3)^2 = 48\,k(n) + 1.
\end{equation}
\end{theorem}

\begin{proof}
Expanding $(4n+3)^2 = 16n^2 + 24n + 9$ and computing
$48\,k(n) + 1 = 48 \cdot \frac{2n^2+3n+1}{6} + 1 = 16n^2 + 24n + 8 + 1 = 16n^2+24n+9$.
\end{proof}

\begin{remark}[Significance]
The identity states that $48\,k(n)+1$ is always a perfect square.
Conversely, suppose $48k+1 = q^2$ for some positive integer $q$.
Then $q$ is necessarily odd (since $48k+1$ is odd) and
$n = (q-3)/4$ is a non-negative integer if and only if
$q \equiv 3 \pmod{4}$. The value $k$ belongs to the sequence
precisely when the additional condition $n \equiv 1, 5 \pmod{6}$ is
satisfied. This perfect-square criterion is the algebraic mechanism
that makes the exact helix construction possible.
\end{remark}

\subsection{The Universal Quantity \texorpdfstring{$Q(k)$}{Q(k)}}

\begin{definition}[Universal quantity]\label{def:Q}
For $k \geq 1$ we define
\begin{equation}\label{eq:Q-def}
  Q(k) \coloneqq \sqrt{48k+1}.
\end{equation}
At the integer values $k = k(n)$, the fundamental identity yields
$Q(k(n)) = 4n+3$.
\end{definition}

\subsection{Exact Inverse Formula}

\begin{theorem}[Inverse of $k(n)$]\label{thm:inverse}
The sequence $k(n)$ possesses the exact inverse
\begin{equation}\label{eq:inverse}
  n(k) = \frac{Q(k) - 3}{4} = \frac{\sqrt{48k+1} - 3}{4}.
\end{equation}
For $k = k(n)$ with $n \equiv 1, 5 \pmod{6}$, this returns the
original index $n$ exactly.
\end{theorem}

\begin{proof}
Solving $(4n+3)^2 = 48k+1$ for $n$:
$4n+3 = \sqrt{48k+1}$, hence $n = (\sqrt{48k+1}-3)/4$.
\end{proof}

\begin{example}
For $k = 46$: $Q = \sqrt{48 \cdot 46 + 1} = \sqrt{2209} = 47$,
hence $n = (47-3)/4 = 11$. Indeed $k(11) = 46$.
\end{example}

\subsection{ODE Characterisation}\label{subsec:ode}

The inverse formula~\eqref{eq:inverse} can also be understood as the
solution of an initial-value problem.

\begin{theorem}[ODE of the height function]\label{thm:ode}
The height function $z = n(k)$ is the unique solution of
\begin{equation}\label{eq:ode}
  \frac{dz}{dk} = \frac{6}{4z+3}, \qquad z\!\left(\tfrac{1}{6}\right) = 0.
\end{equation}
The explicit solution is $2z^2 + 3z + 1 = 6k$, which upon solving for
$z \geq 0$ coincides with~\eqref{eq:inverse}.
\end{theorem}

\begin{proof}
Separation of variables: $(4z+3)\,dz = 6\,dk$.
Integration: $2z^2 + 3z = 6k + C$.
The initial condition $z(1/6) = 0$ gives $C = -1$, hence
$2z^2+3z+1 = 6k$, i.e.\ $k(z) = k$.
\end{proof}

\begin{remark}[Asymptotics]
The slope $dz/dk = 6/(4z+3) \approx 3/(2z)$ for large $z$ yields
$z \approx \sqrt{3k}$. This is the leading term of the exact formula
$z = (\sqrt{48k+1}-3)/4 \approx \sqrt{3k} - 3/4 + O(k^{-1/2})$.
\end{remark}

\subsection{All Helix Quantities as Functions of \texorpdfstring{$Q$}{Q}}
\label{subsec:Q-table}

The quantity $Q$ determines all geometric coordinates of the helix
through exact, closed-form expressions:

\begin{table}[H]
\centering
\renewcommand{\arraystretch}{1.6}
\caption{Helix coordinates as exact functions of $Q(k) = \sqrt{48k+1}$.
  At integer values $k = k(n)$, one has $Q = 4n+3$.}
\label{tab:Q-overview}
\begin{tabular}{lcc}
\toprule
Quantity & Formula in $Q$ & Value at $Q = 4n+3$ \\
\midrule
Height $z$
  & $\dfrac{Q-3}{4}$
  & $n$ \\[8pt]
Angle $\theta$
  & $\dfrac{\pi\,Q}{12}$
  & $\dfrac{\pi}{3}n + \dfrac{\pi}{4}$ \\[8pt]
Radius $r$
  & $\dfrac{Q}{2\pi}$
  & $\dfrac{4n+3}{2\pi}$ \\[8pt]
Function value $k$
  & $\dfrac{Q^2-1}{48}$
  & $\dfrac{2n^2+3n+1}{6}$ \\[8pt]
Derivative $k'(n)$
  & $\dfrac{Q}{6}$
  & $\dfrac{4n+3}{6}$ \\
\bottomrule
\end{tabular}
\end{table}

The table reveals that $Q$ is a \emph{universal parameter}: a single
square root encodes the complete three-dimensional geometry of the helix.


% ── Chapter 4: The Conical Helix ─────────────────────────────
% ============================================================
%  Chapter 4 — The Conical Helix: Direct Construction
% ============================================================
\section{The Conical Helix}\label{sec:helix}

Using the universal quantity $Q(k) = \sqrt{48k+1}$ from
Section~\ref{sec:identity}, we now construct the conical helix
directly and prove its geometric properties as algebraic corollaries.

\subsection{Exact Parametrisation}

\begin{theorem}[Exact helix parametrisation]\label{thm:helix}
Let $Q(k) = \sqrt{48k+1}$. The conical helix is given exactly by
\begin{equation}\label{eq:helix-exact}
  \boxed{
  \vec{p}(k) =
  \begin{pmatrix}
    \dfrac{Q}{2\pi}\cos\!\Bigl(\dfrac{\pi\,Q}{12}\Bigr)\\[10pt]
    \dfrac{Q}{2\pi}\sin\!\Bigl(\dfrac{\pi\,Q}{12}\Bigr)\\[10pt]
    \dfrac{Q-3}{4}
  \end{pmatrix},
  \qquad Q = Q(k).
  }
\end{equation}
At the integer values $k = k(n)$ with $n \equiv 1, 5 \pmod{6}$,
one has $Q = 4n+3$, and the point $\vec{p}$ has integer height $z = n$.
\end{theorem}

\begin{proof}
By Definition~\ref{def:Q} and Table~\ref{tab:Q-overview}, the three
coordinates are $r = Q/(2\pi)$, $\theta = \pi Q/12$, and $z = (Q-3)/4$.
At $k = k(n)$ the fundamental identity (Theorem~\ref{thm:fundamental})
gives $Q = 4n+3$, whence $z = (4n+3-3)/4 = n$ and
$\theta = \pi(4n+3)/12 = \pi n/3 + \pi/4$.
\end{proof}

Equivalently, the helix admits the standard $\theta$-parametrisation:

\begin{corollary}[$\theta$-parametrisation]\label{cor:theta-param}
With $a = 6/\pi^2$ and $c_1 = 3/\pi$:
\begin{equation}\label{eq:helix-theta}
  \gamma(\theta) =
  \begin{pmatrix}
    a\,\theta\cos\theta \\[4pt]
    a\,\theta\sin\theta \\[4pt]
    -\tfrac{3}{4} + c_1\,\theta
  \end{pmatrix}.
\end{equation}
The helix intersects the plane $z = n$ at the angle
$\theta_n = \frac{\pi}{3}n + \frac{\pi}{4}$.
\end{corollary}

\begin{proof}
Setting $\theta = \pi Q/12$ and $r = Q/(2\pi)$ gives
$r = (6/\pi^2)\theta = a\theta$. The height is
$z = (Q-3)/4 = (12\theta/\pi - 3)/4 = (3/\pi)\theta - 3/4 = c_1\theta - 3/4$.
\end{proof}

\subsection{Spiral Conformity}\label{subsec:spiral}

\begin{theorem}[Archimedean spiral]\label{thm:spiral}
The helix points $(r(k), \theta(k))$ lie exactly on the Archimedean
spiral $r = \frac{6}{\pi^2}\,\theta$.
\end{theorem}

\begin{proof}
Direct verification:
$\frac{6}{\pi^2}\cdot\frac{\pi Q}{12}
  = \frac{Q}{2\pi}
  = r(k)$.
\end{proof}

\begin{remark}[Connection to the Basel problem]
The spiral slope constant $a = 6/\pi^2 = 1/\zeta(2)$ is the reciprocal
of the Riemann zeta function at $s = 2$, the solution of Euler's
Basel problem~\cite{Euler1740, Hardy2008}.
\end{remark}

\subsection{Cone Conformity}\label{subsec:cone}

\begin{theorem}[Cone equation]\label{thm:cone}
The helix points $(r(k), z(k))$ lie exactly on the cone of revolution
$z = -\frac{3}{4} + \frac{\pi}{2}\,r$.
\end{theorem}

\begin{proof}
Direct verification:
$-\frac{3}{4} + \frac{\pi}{2}\cdot\frac{Q}{2\pi}
  = -\frac{3}{4} + \frac{Q}{4}
  = \frac{Q-3}{4}
  = z(k)$.
\end{proof}

\begin{remark}[Geometric interpretation]
The cone has apex at $(0, 0, -3/4)$, half-angle
$\alpha = \arctan(2/\pi) \approx 32.48°$, and the helix rises by
exactly $\Delta z = 6$ per full revolution
($\Delta\theta = 2\pi$, $\Delta z = c_1 \cdot 2\pi = 6$).
\end{remark}

\subsection{Three Invariants from One Identity}

\begin{tcolorbox}[
  title={Summary: Three simultaneous invariants},
  colback=blue!3,
  colframe=blue!40!black,
  enhanced
]
The parametrisation~\eqref{eq:helix-exact} satisfies three exact
invariants, all algebraic consequences of the single identity
$(4n+3)^2 = 48k(n)+1$:
\begin{enumerate}[label=(\roman*)]
  \item \textbf{Arithmetic invariant:} $Q(k(n)) = 4n+3$
        \quad (Theorem~\ref{thm:fundamental}),
  \item \textbf{Spiral invariant:} $r = \frac{6}{\pi^2}\,\theta$
        \quad (Theorem~\ref{thm:spiral}),
  \item \textbf{Cone invariant:} $z = -\frac{3}{4} + \frac{\pi}{2}\,r$
        \quad (Theorem~\ref{thm:cone}).
\end{enumerate}
\end{tcolorbox}

\subsection{Velocity, Curvature, and Torsion}\label{subsec:curvature}

The speed of the helix in the $\theta$-parametrisation is obtained by
differentiating~\eqref{eq:helix-theta}. The cross terms cancel via the
Pythagorean identity, yielding:
\begin{equation}\label{eq:speed}
  \|\gamma'(\theta)\|
  = \frac{3}{\pi^2}\sqrt{4(1+\theta^2) + \pi^2}.
\end{equation}

\begin{theorem}[Exact 3D curvature and torsion]\label{thm:curvature-torsion}
With $a = 6/\pi^2$ and $c_1 = 3/\pi$:
\begin{align}
  \kappa(\theta)
  &= \frac{a\,\sqrt{c_1^2(\theta^2+4) + a^2(\theta^2+2)^2}}
          {\bigl(a^2(1+\theta^2) + c_1^2\bigr)^{3/2}},
  \label{eq:kappa}\\[6pt]
  \tau(\theta)
  &= \frac{c_1\,(\theta^2 + 6)}
          {c_1^2(\theta^2+4) + a^2(\theta^2+2)^2}.
  \label{eq:tau}
\end{align}
\end{theorem}

\begin{proof}
Computing the derivatives $\gamma'$, $\gamma''$, $\gamma'''$ of the
parametrisation~\eqref{eq:helix-theta}, one obtains the cross product
\[
  \gamma' \times \gamma''
  = \bigl(-ac_1(2\cos\theta - \theta\sin\theta),\;
          -ac_1(2\sin\theta + \theta\cos\theta),\;
          a^2(\theta^2+2)\bigr)
\]
with $|\gamma' \times \gamma''|^2 = a^2[c_1^2(\theta^2+4) +
a^2(\theta^2+2)^2]$, $|\gamma'|^2 = a^2(1+\theta^2) + c_1^2$, and
the scalar triple product
$(\gamma' \times \gamma'') \cdot \gamma''' = a^2 c_1(\theta^2+6)$.
The formulae follow from $\kappa = |\gamma' \times \gamma''|/|\gamma'|^3$
and $\tau = (\gamma' \times \gamma'') \cdot \gamma''' / |\gamma' \times
\gamma''|^2$.
\end{proof}

\begin{remark}[Asymptotics]
For $\theta \to \infty$:
$\kappa(\theta) \sim 1/(a\theta) \to 0$ (the helix straightens) and
$\tau(\theta) \to c_1/a^2 = \pi^3/12$ (the torsion converges to a
non-zero constant).
\end{remark}

\subsection{Angular Spacing Pattern}\label{subsec:angles}

\begin{proposition}[Alternating angular spacings]\label{prop:angles}
Between consecutive integer points the angular spacings are:
\begin{align}
  \Delta\theta^{(4)} &= \theta_{n+4} - \theta_n = \tfrac{4\pi}{3}
  \quad\text{(long step)},
  \label{eq:angle-long}\\
  \Delta\theta^{(2)} &= \theta_{n+2} - \theta_n = \tfrac{2\pi}{3}
  \quad\text{(short step)}.
  \label{eq:angle-short}
\end{align}
Together, $\Delta\theta^{(4)} + \Delta\theta^{(2)} = 2\pi$,
so that each long--short pair completes exactly one full revolution.
\end{proposition}

\begin{proof}
From $\theta_n = \frac{\pi}{3}n + \frac{\pi}{4}$:
$\theta_{n+\Delta n} - \theta_n = \frac{\pi}{3}\Delta n$.
\end{proof}


% ── Chapter 5: Arc Length and Asymptotic Equivalence ──────────
% ============================================================
%  Chapter 5 — Arc Length and Asymptotic Equivalence
% ============================================================
\section{Arc Length and Asymptotic Equivalence}\label{sec:arclength}

In this section we derive the arc-length formula for the parabola $k(n)$
and prove that the arc length converges asymptotically to the
difference values $\Delta k$. The universal quantity $Q$ provides a
unified framework for comparing the arc lengths of the parabola and the
helix.

\subsection{Arc-Length Formula for the Parabola}

The arc-length element of the curve $(n, k(n))$ is
$ds = \sqrt{1 + [k'(n)]^2}\,dn$ with $k'(n) = (4n+3)/6$.
Substituting $u = 4n+3$, $du = 4\,dn$:

\begin{theorem}[Closed-form arc length]\label{thm:arc-length-parabola}
Let $u_i = 4n_i + 3$. The arc length of $k(n)$ between $n_1$ and $n_2$ is:
\begin{equation}\label{eq:arc-length-closed}
  L(n_1, n_2)
  = \frac{1}{48}\Bigl[u\sqrt{u^2+36}
    + 36\operatorname{arcsinh}\!\left(\frac{u}{6}\right)
    \Bigr]_{u_1}^{u_2}.
\end{equation}
\end{theorem}

\begin{proof}
From $ds = \frac{1}{6}\sqrt{36 + (4n+3)^2}\,dn$ and $u = 4n+3$:
$L = \frac{1}{24}\int_{u_1}^{u_2}\sqrt{u^2+36}\,du$.
Applying the standard integral
$\int\sqrt{u^2+a^2}\,du = \frac{1}{2}[u\sqrt{u^2+a^2} +
a^2\operatorname{arcsinh}(u/a)]$ with $a = 6$ and combining the
prefactors $\frac{1}{24}\cdot\frac{1}{2} = \frac{1}{48}$ yields the
result.
\end{proof}

\subsection{Unified \texorpdfstring{$Q$}{Q}-Representation}
\label{subsec:Q-arclength}

Since $u = 4n+3 = Q(k(n))$, the arc-length element can be written
directly in terms of $Q$:

\begin{theorem}[$Q$-parametrisation of the arc lengths]\label{thm:Q-arc}
The arc-length elements of the parabola and of the conical helix
admit the unified form
\begin{equation}\label{eq:ds-Q}
  ds = \frac{1}{24}\sqrt{Q^2 + C}\; dQ,
\end{equation}
where
\[
  C_{\mathrm{Parabola}} = 36,
  \qquad
  C_{\mathrm{Helix}} = 36 + \frac{144}{\pi^2} \approx 50.59.
\]
\end{theorem}

\begin{proof}
For the parabola, this is immediate from the substitution $u = Q$.
For the helix, using $\theta = \pi Q/12$ and $d\theta = (\pi/12)\,dQ$
in the speed formula~\eqref{eq:speed}:
\begin{align*}
  ds &= \|\gamma'(\theta)\|\,d\theta
     = \frac{3}{\pi^2}\sqrt{4(1+\theta^2)+\pi^2}\cdot\frac{\pi}{12}\,dQ \\
     &= \frac{1}{4\pi}\sqrt{4+4\theta^2+\pi^2}\,dQ
     = \frac{1}{24}\sqrt{Q^2 + 36 + \tfrac{144}{\pi^2}}\,dQ.
     \qedhere
\end{align*}
\end{proof}

\begin{remark}[Convergence of the arc lengths]
The difference $C_{\mathrm{Helix}} - C_{\mathrm{Parabola}} = 144/\pi^2$
is the exact contribution of the spiral motion. For $Q \to \infty$ the
ratio
$ds_{\mathrm{Helix}}/ds_{\mathrm{Parabola}} = \sqrt{(Q^2+C_H)/(Q^2+C_P)}
\to 1$, so the arc lengths of both curves converge.
\end{remark}

\subsection{Asymptotic Equivalence: Arc Length \texorpdfstring{$\approx \Delta k$}{≈ Δk}}

The most striking consequence of the arc-length formula is that the
difference values $\Delta k$ of the sequence are themselves excellent
approximations of the arc length.

\begin{proposition}[Asymptotic arc-length--difference equivalence]
\label{prop:asymptotic}
Let $L^{(2)}(m)$ and $L^{(4)}(m)$ denote the arc lengths of the
parabola over the $m$-th short step and the $m$-th long step,
respectively (Definition~\ref{def:steps}). Then, as $m \to \infty$:
\[
  \frac{L^{(2)}(m)}{\delta^{(2)}_m} \to 1,
  \qquad
  \frac{L^{(4)}(m)}{\delta^{(4)}_m} \to 1.
\]
\end{proposition}

\begin{proof}
We demonstrate the short-step case; the long-step case is analogous.
The short step connects $n_1 = 6m-1$ to $n_2 = 6m+1$, giving
$u_1 = 24m-1$ and $u_2 = 24m+7$.

By Theorem~\ref{thm:arc-length-parabola}, the arc length is
$L^{(2)}(m) = \frac{1}{48}[F(u)]_{u_1}^{u_2}$ where
$F(u) = u\sqrt{u^2+36} + 36\operatorname{arcsinh}(u/6)$.
For large $u$ the expansion
$\sqrt{u^2+36} = u + 18/u + O(u^{-3})$ gives:
\[
  F(u) = u^2 + 18 + 36\ln u + O(u^{-1}).
\]
Hence
\[
  F(u_2) - F(u_1)
  = (u_2^2 - u_1^2) + 36\ln\frac{u_2}{u_1} + O(m^{-1}).
\]
Now $u_2^2 - u_1^2 = (u_2+u_1)(u_2-u_1) = (48m+6)\cdot 8 = 384m + 48$,
so
\[
  L^{(2)}(m) = \frac{384m+48}{48} + O\!\left(\frac{\ln m}{m}\right)
             = 8m + 1 + O\!\left(\frac{\ln m}{m}\right)
             = \delta^{(2)}_m + o(1).
             \qedhere
\]
\end{proof}

\begin{table}[H]
\centering
\caption{Arc length vs.\ difference value. The ratio $L/\Delta k$
  converges monotonically to~$1$.}
\renewcommand{\arraystretch}{1.15}
\begin{tabular}{cccccc}
\toprule
$n_1 \to n_2$ & $\Delta n$ & $\Delta k$ & Arc length $L$ & $L / \Delta k$ \\
\midrule
 $1 \to  5$ & 4 & 10 & 10.839 & 1.0839 \\
 $5 \to  7$ & 2 &  9 &  9.221 & 1.0246 \\
 $7 \to 11$ & 4 & 26 & 26.310 & 1.0119 \\
$11 \to 13$ & 2 & 17 & 17.117 & 1.0069 \\
$13 \to 17$ & 4 & 42 & 42.191 & 1.0045 \\
$17 \to 19$ & 2 & 25 & 25.080 & 1.0032 \\
$19 \to 23$ & 4 & 58 & 58.138 & 1.0024 \\
$23 \to 25$ & 2 & 33 & 33.061 & 1.0018 \\
\bottomrule
\end{tabular}
\end{table}

\begin{remark}[Geometric meaning]
This convergence has a transparent cause: the slope
$k'(n) = (4n+3)/6 = Q/6$ grows without bound, so the horizontal
component $\Delta n \leq 4$ contributes ever less to the arc-length
element $\sqrt{1 + (k')^2}$. In the limit, $ds/dn \to k'(n)$ and
hence $L \to \int k'\,dn = \Delta k$.
\end{remark}

\subsection{Taylor Approximation}

The asymptotic equivalence $L \approx \Delta k$ can also be confirmed
through the $p$-parametrisation of the \emph{spiral} arc length
(cf.\ Section~\ref{sec:helix}), where the integrand takes the form
$\sqrt{1 + 3/(\pi^2 p)}$. Expanding for large $p$:

\begin{theorem}[Taylor approximation of the spiral arc length]
\label{thm:taylor}
For the spiral arc length in the $p$-parametrisation, with
$p_i = k(n_i)$:
\begin{equation}\label{eq:taylor}
  L \approx (p_2 - p_1)
  + \frac{3}{2\pi^2}\ln\!\left(\frac{p_2}{p_1}\right)
  + \frac{9}{8\pi^4}\left(\frac{1}{p_1} - \frac{1}{p_2}\right).
\end{equation}
The leading term $(p_2 - p_1) = \Delta k$ confirms the asymptotic
equivalence; the logarithmic correction is of relative order
$O(\ln p / p)$.
\end{theorem}


% ── Chapter 6: Prime Numbers on the Helix ────────────────────
% ============================================================
%  Chapter 6 — Prime Numbers on the Helix
% ============================================================
\section{Prime Numbers on the Helix}\label{sec:primes}

The conical helix provides a natural geometric embedding of the
congruence classes $n \equiv 1, 5 \pmod{6}$. Since these are exactly the
classes containing all primes greater than~$3$
(Theorem~\ref{thm:primes-mod6}), the primes appear as a distinguished
subset of the integer helix points.

\subsection{Integer Points and the Prime Corollary}

By Corollary~\ref{cor:theta-param}, the helix intersects the plane
$z = n$ at the angle $\theta_n = \frac{\pi}{3}n + \frac{\pi}{4}$
for all $n \equiv 1, 5 \pmod{6}$. Combining this with the classical
fact that all primes $p > 3$ lie in these congruence classes yields
the central geometric corollary:

\begin{theorem}[Primes on the helix]\label{thm:primes-helix}
All primes $p > 3$ lie as discrete points on the conical helix at angles
\begin{equation}\label{eq:prime-angle}
  \boxed{\theta_p = \frac{\pi}{3}\,p + \frac{\pi}{4}.}
\end{equation}
\end{theorem}

\begin{proof}
\textbf{Step~1.} Every prime $p > 3$ satisfies
$p \equiv 1$ or $5 \pmod{6}$ (Theorem~\ref{thm:primes-mod6}).

\textbf{Step~2.} The helix attains integer height $z = n$ precisely at
$\theta_n = \frac{\pi}{3}n + \frac{\pi}{4}$ for
$n \equiv 1, 5 \pmod{6}$ (Corollary~\ref{cor:theta-param}).

\textbf{Conclusion.} Since every prime $p > 3$ satisfies the condition
from Step~1, it lies on the helix at $\theta_p$.
\end{proof}

\begin{example}[First primes on the helix]
\begin{center}
\begin{tabular}{cccc}
\toprule
$p$ & $\theta_p$ (exact) & $\theta_p$ (numerical) & $k(p)$ \\
\midrule
 5 & $23\pi/12$ &  6.021 &  11 \\
 7 & $31\pi/12$ &  8.116 &  20 \\
11 & $47\pi/12$ & 12.305 &  46 \\
13 & $55\pi/12$ & 14.399 &  63 \\
17 & $71\pi/12$ & 18.588 & 105 \\
19 & $79\pi/12$ & 20.682 & 130 \\
23 & $95\pi/12$ & 24.871 & 188 \\
\bottomrule
\end{tabular}
\end{center}
\end{example}

\subsection{Composite Numbers and Complete Characterisation}

The helix does not contain only primes. Since it embeds all indices
$n \equiv 1, 5 \pmod{6}$, composite numbers from these classes also
appear. The first composite on the helix is $n = 25 = 5^2$.

\begin{remark}[Complete characterisation]
The helix contains exactly the set
$\{n \in \mathbb{Z}_{>0} : n \equiv 1, 5 \pmod{6}\}$.
The primes $p > 3$ form a distinguished but non-isolable subset.
The numbers $2$, $3$, and all $n$ with $n \equiv 0, 2, 3, 4 \pmod{6}$
do not lie on the helix.
\end{remark}

\subsection{Angular Spacings and Prime Gaps}

For consecutive primes $p < q$ with $p, q > 3$, the angular gap
$\theta_q - \theta_p = \frac{\pi}{3}(q-p)$ depends only on the
prime gap $g = q - p$. Two particularly noteworthy cases:

\begin{center}
\begin{tabular}{lcc}
\toprule
Prime type & Gap $g$ & Angular spacing \\
\midrule
Twin primes ($q = p+2$) & 2 & $2\pi/3$ \\
Cousin primes ($q = p+4$) & 4 & $4\pi/3$ \\
Sexy primes ($q = p+6$) & 6 & $2\pi$ (full revolution) \\
\bottomrule
\end{tabular}
\end{center}

Twin primes thus occupy adjacent positions on the same $120°$-arc,
while sexy primes are vertically aligned on the helix.

\subsection{Asymptotic Density}

\begin{proposition}[Density of primes on the helix]\label{prop:density}
The density of primes among the integer helix points tends to zero:
\[
  \frac{\#\{p \leq N : p \text{ prime},\; p > 3\}}
       {\#\{n \leq N : n \equiv 1, 5 \pmod{6}\}}
  \;\sim\; \frac{N/\ln N}{N/3}
  \;=\; \frac{3}{\ln N}
  \;\xrightarrow{N\to\infty}\; 0.
\]
\end{proposition}

\subsection{Comparison with Other Prime Spirals}

\begin{remark}[Derived vs.\ empirical]
In contrast to the Ulam spiral and the Sacks spiral, which represent
empirically discovered patterns, the present helix is \emph{derived}
algebraically from the quadratic sequence $k(n)$. The appearance of all
primes $p > 3$ is a \emph{theorem} (Theorem~\ref{thm:primes-helix}),
not an observation. The underlying mechanism --- that the congruence
classes $1, 5 \pmod{6}$ are simultaneously the integrality locus of
$k(n)$ and the residue classes coprime to~$6$ --- is an algebraic
coincidence rooted in the factorisation
$2n^2+3n+1 = (2n+1)(n+1)$.
\end{remark}

\subsection{Open Questions}

\begin{enumerate}
  \item Can the helix structure yield insights into prime gaps or the
    distribution of twin primes?
  \item Does the arc-length parametrisation via $Q(k)$ provide a
    natural metric on the primes?
  \item Which other quadratic sequences lead to analogous helix
    constructions, and what are their geometric invariants?
\end{enumerate}


% ── Bibliography ─────────────────────────────────────────────
\printbibliography

\end{document}
